
\section{Definitions}

Now, as we are done talking about the basics of representation theory, we get to the main goal of this exposition: introducing Semisimple algebras. The main claim we aim to prove in this chapter is that all semisimple algebras are product of matrix algebras. This characterization shows that semisimple algebras are the easiest to work with.

\begin{definition}[Semisimple Algebra]
    An algebra $A$, viewed as a representation of $A$ is \emph{semisimple} or \emph{completely reducible} if it is a direct sum of irreducible representations of $A$.
    
\end{definition}
There are many ways to define semisimple algebras, and they are all equivalent. At the end, we see that all of these definitions are just stating properties that such an algebra has.
We will prove that all of these properties are equivalent.

In fact, semisimple algebras are nice to work with since they are all isomorphic to a direct sum of matrix algebra (Artin-Wedderburn). Moreover, we will see a characterization for an arbitrary group algebra over a field (with some restrictions on the characterstic of that field) to be semisimple thus nice, and this is the main title of this exposition which is Maschke's Theorem.

Now, we have a nice lemma about irreducible representations that makes a lot of sense if we can see that the kernel and the image of an intertwining operator are subrepresentations.

\begin{lemma}
    Let $A$ be a $k-$algebra and $V,W$ be two irreducible representations of $A$. Let $\varphi: V \to W$ be an intertwining operator. Then $\varphi$ is either zero or an isomorphism. 
\end{lemma}
\begin{proof}
    Since $\ker \varphi$ is a subrepresentation of $V$ then $\ker V ={0}$ or $\ker V = V$. 
    \begin{enumerate}
        \item Case 1: If $\ker V =0 $, then $\varphi$ is injective. But also $W$ is irreducible so $\operatorname{Img} \varphi = {0}$ or $\operatorname{Img} \varphi = W$ by rank-nullity we get that $\operatorname{Img} \varphi $ must be $W$ hence $\varphi$ is an isomorphism.
        \item Case 2: if $\ker V = V$ then $\varphi$ is trivially the zero map. 
    \end{enumerate}
    and this ends the proof.
\end{proof}

before we go to the main theorems we define the regular and trivial representations

\begin{definition}[Regular Representations]

\end{definition}

\begin{definition}[Trivial Representation]
    
\end{definition}

\begin{definition}[Radical]
    
\end{definition}
\begin{theorem}[Finitely many irreducible representations]
    
\end{theorem}

\begin{theorem}[The irreducible representations of a semisimple algebra]
    
\end{theorem}
\begin{theorem}[Equivalences to semisimplicity]
    

\end{theorem}



\begin{theorem}[Density Theorem]
    
\end{theorem}

\begin{theorem}[Maschke Theorem]
    
\end{theorem}